% GENERATED FILE. Edit canonical lessons, then run npm run generate:books.
% canonical-source-hash: fnv1a64:e0b44616619f88a6
% canonical-lessons: PT-C426-acampar, PT-C426-caravana, PT-C426-autocaravana, PT-C426-estaca, PT-C426-mosquiteiro

\chapter{No parque de campismo}
\label{ch:pt-no-parque-de-campismo}
\hlchaptermodality{426}

\begin{chapteropening}
\textbf{By the end of this chapter:} \emph{I can talk about camping: camping out, the caravan, the camper van, the tent peg and the mosquito net.}

Complete the last lesson of chapter 426: I can talk about camping: camping out, the caravan, the camper van, the tent peg and the mosquito net.
\end{chapteropening}

\section[acampar]{acampar — to camp}

\label{lesson:PT-C426-acampar}



\begin{quote}
Before the new one: say the Portuguese for the guesthouse, then the Portuguese for camping.
\end{quote}

\subsection*{You'll want to know: acampar}
\textbf{acampar} — \textquotedblleft{}to camp\textquotedblright{}.

\begin{quote}
Acampámos perto de um lago.
\end{quote}

\begin{quote}
\emph{We camped near a lake.}
\end{quote}

\emph{Acampar} in the wild is not allowed everywhere.

\begin{cousinweb}
\emph{acampar} was borrowed from Italian \emph{accampare} 'to set up camp' (EtymDB), built on \emph{campo} 'field', from Latin \emph{campus}.
\end{cousinweb}

\begin{grammarlens}[title={What you've built}]
One more word for telling, explaining and describing at B1.
\end{grammarlens}

\subsection*{Your turn}
\begin{itemize}
\raggedright
  \item \emph{Say it:} \emph{acampar}
  \item \emph{Say it:} \emph{acampar}, once more
  \item \emph{Say it:} a sentence of your own with \emph{acampar}
  \item \emph{Recall:} say the Portuguese for the guesthouse, then the Portuguese for camping, then say \emph{acampar} again
\end{itemize}

\begin{tcolorbox}[breakable,colback=teal!4,colframe=teal!35!black,title={Before you move on}]
What does acampar mean? (\textquotedblleft{}To camp\textquotedblright{}.) Say it once more.
\end{tcolorbox}
\section[a caravana]{a caravana — the caravan}

\label{lesson:PT-C426-caravana}



\begin{quote}
Before the new one: say the Portuguese for camping, then the Portuguese for to camp.
\end{quote}

\subsection*{You'll want to know: a caravana}
\textbf{a caravana} — \textquotedblleft{}the caravan\textquotedblright{}.

\begin{quote}
Os vizinhos passam o verão numa caravana.
\end{quote}

\begin{quote}
\emph{The neighbours spend the summer in a caravan.}
\end{quote}

Also a line of travellers, as in the desert.

\begin{cousinweb}
\emph{caravana} comes from Persian \emph{kārvān} 'group of travellers' (EtymDB).
\end{cousinweb}

\begin{grammarlens}[title={What you've built}]
One more word for telling, explaining and describing at B1.
\end{grammarlens}

\subsection*{Your turn}
\begin{itemize}
\raggedright
  \item \emph{Say it:} \emph{a caravana}
  \item \emph{Say it:} \emph{a caravana}, once more
  \item \emph{Say it:} a sentence of your own with \emph{a caravana}
  \item \emph{Recall:} say the Portuguese for camping, then the Portuguese for to camp, then say \emph{a caravana} again
\end{itemize}

\begin{tcolorbox}[breakable,colback=teal!4,colframe=teal!35!black,title={Before you move on}]
What does a caravana mean? (\textquotedblleft{}The caravan\textquotedblright{}.) Say it once more.
\end{tcolorbox}
\section[a autocaravana]{a autocaravana — the camper van, the motorhome}

\label{lesson:PT-C426-autocaravana}



\begin{quote}
Before the new one: say the Portuguese for to camp, then the Portuguese for the caravan.
\end{quote}

\subsection*{You'll want to know: a autocaravana}
\textbf{a autocaravana} — \textquotedblleft{}the camper van, the motorhome\textquotedblright{}.

\begin{quote}
Viajámos pelo país de autocaravana.
\end{quote}

\begin{quote}
\emph{We travelled round the country in a camper van.}
\end{quote}

In Brazil: \emph{o motorhome}.

\begin{cousinweb}
\emph{autocaravana} is built from \emph{auto-} 'self-moving', from Greek \emph{autós} 'self', and \emph{caravana}, which comes from Persian \emph{kārvān} (EtymDB).
\end{cousinweb}

\begin{grammarlens}[title={What you've built}]
One more word for telling, explaining and describing at B1.
\end{grammarlens}

\subsection*{Your turn}
\begin{itemize}
\raggedright
  \item \emph{Say it:} \emph{a autocaravana}
  \item \emph{Say it:} \emph{a autocaravana}, once more
  \item \emph{Say it:} a sentence of your own with \emph{a autocaravana}
  \item \emph{Recall:} say the Portuguese for to camp, then the Portuguese for the caravan, then say \emph{a autocaravana} again
\end{itemize}

\begin{tcolorbox}[breakable,colback=teal!4,colframe=teal!35!black,title={Before you move on}]
What does a autocaravana mean? (\textquotedblleft{}The camper van, the motorhome\textquotedblright{}.) Say it once more.
\end{tcolorbox}
\section[a estaca]{a estaca — the tent peg; the stake}

\label{lesson:PT-C426-estaca}



\begin{quote}
Before the new one: say the Portuguese for the caravan, then the Portuguese for the camper van.
\end{quote}

\subsection*{You'll want to know: a estaca}
\textbf{a estaca} — \textquotedblleft{}the tent peg; the stake\textquotedblright{}.

\begin{quote}
O vento arrancou as estacas da tenda.
\end{quote}

\begin{quote}
\emph{The wind pulled out the tent pegs.}
\end{quote}

A thicker \emph{estaca} holds up a young tree or a fence.

\begin{cousinweb}
\emph{estaca} is usually traced to a Germanic word for a pointed stick; EtymDB lists English \emph{stake} and Old French \emph{estaque} as its cousins.
\end{cousinweb}

\begin{grammarlens}[title={What you've built}]
One more word for telling, explaining and describing at B1.
\end{grammarlens}

\subsection*{Your turn}
\begin{itemize}
\raggedright
  \item \emph{Say it:} \emph{a estaca}
  \item \emph{Say it:} \emph{a estaca}, once more
  \item \emph{Say it:} a sentence of your own with \emph{a estaca}
  \item \emph{Recall:} say the Portuguese for the caravan, then the Portuguese for the camper van, then say \emph{a estaca} again
\end{itemize}

\begin{tcolorbox}[breakable,colback=teal!4,colframe=teal!35!black,title={Before you move on}]
What does a estaca mean? (\textquotedblleft{}The tent peg; the stake\textquotedblright{}.) Say it once more.
\end{tcolorbox}
\section[o mosquiteiro]{o mosquiteiro — the mosquito net}

\label{lesson:PT-C426-mosquiteiro}



\begin{quote}
Before the new one: say the Portuguese for the camper van, then the Portuguese for the tent peg.
\end{quote}

\subsection*{You'll want to know: o mosquiteiro}
\textbf{o mosquiteiro} — \textquotedblleft{}the mosquito net\textquotedblright{}.

\begin{quote}
Sem mosquiteiro, ninguém dorme no verão.
\end{quote}

\begin{quote}
\emph{Without a mosquito net, nobody sleeps in summer.}
\end{quote}

Built on \emph{o mosquito}.

\begin{cousinweb}
\emph{mosquiteiro} is built on \emph{mosquito}, a small form of \emph{mosca} 'fly', which is inherited from Latin \emph{musca} (Lewis \& Short).
\end{cousinweb}

\begin{grammarlens}[title={What you've built}]
One more word for telling, explaining and describing at B1.
\end{grammarlens}

\subsection*{Your turn}
\begin{itemize}
\raggedright
  \item \emph{Say it:} \emph{o mosquiteiro}
  \item \emph{Say it:} \emph{o mosquiteiro}, once more
  \item \emph{Say it:} a sentence of your own with \emph{o mosquiteiro}
  \item \emph{Recall:} say the Portuguese for the camper van, then the Portuguese for the tent peg, then say \emph{o mosquiteiro} again
\end{itemize}

\begin{tcolorbox}[breakable,colback=teal!4,colframe=teal!35!black,title={Before you move on}]
What does o mosquiteiro mean? (\textquotedblleft{}The mosquito net\textquotedblright{}.) Say it once more.
\end{tcolorbox}
